\documentclass[11pt]{amsart}
\usepackage[margin=1in]{geometry}
\usepackage{amsmath,amssymb,amsthm}
\usepackage{longtable,booktabs,array}
\usepackage{url}
\usepackage[hidelinks]{hyperref}
\newtheorem{theorem}{Theorem}[section]
\newtheorem{proposition}[theorem]{Proposition}
\newtheorem{lemma}[theorem]{Lemma}
\newtheorem{corollary}[theorem]{Corollary}
\theoremstyle{definition}
\newtheorem{definition}[theorem]{Definition}
\theoremstyle{remark}
\newtheorem{remark}[theorem]{Remark}
\newcommand{\lean}[1]{\path{#1}}
\sloppy

\title[Elementary certified facts for the Fermat quintic fourfold]{Elementary Certified Facts for the Fermat Quintic Fourfold $V(F_5)\subset\mathbb{P}^5$: the Euler Identity, an Euler Defect, and the Isolated Integer Critical Point of the Affine Cone}
\author{Benjamin Stanley Frohman}
\email{frohmanbenjamin@gmail.com}
\date{November 4, 2026}
\begin{document}
\begin{abstract}
We formalize in Lean~4, using the core library only, a small set of elementary facts about the Fermat quintic polynomial $F_5=\sum_{i=0}^5 x_i^5$ evaluated on integers: Euler's identity $5F_5=\sum x_i\,\partial_iF_5$; the exact Euler defect $-5\prod x_i$ of the non-homogeneous perturbation $G=F_5-5\prod x_i$; and the fact that the only integer point at which all six partial derivatives $5x_i^4$ vanish is the origin. We also record, as checked arithmetic and string labels, the bookkeeping constants $126-6=120$ and $1751$ that appear in Hodge-theoretic discussions of this host, and we state explicitly that no cohomology group, Hodge number, or algebraic cycle is constructed. All proofs are free of \texttt{sorry} and use only Lean's standard axioms.
\end{abstract}
\maketitle

\section{Introduction}
The Fermat quintic fourfold $V(F_5)=\{x_0^5+\dots+x_5^5=0\}\subset\mathbb{P}^5$ is a standard smooth test case for Hodge-theoretic computations. Smoothness of a hypersurface $V(F)$ of degree $d$ in characteristic zero is the statement that the partial derivatives of $F$ have no common zero other than the origin of the affine cone, and Euler's identity $dF=\sum x_i\partial_iF$ shows that such a common zero automatically lies on $V(F)$.

This paper records the parts of this picture that are formalized, in the precise form they are formalized: polynomial identities and an implication about integer points. Everything is stated over $\mathbb{Z}$ with the derivatives entered as explicit polynomials. The repository's \lean{HostLabel} record is a documentation label with string fields; it is not a scheme.

\section{Definitions}
For $x_0,\dots,x_5\in\mathbb{Z}$ define
\[
F_5=\sum_{i=0}^5 x_i^5,\qquad D_i=5x_i^4\ (i=0,\dots,5),\qquad G=F_5-5\,x_0x_1x_2x_3x_4x_5 .
\]
In the formalization, $D_i$ is \emph{defined} to be $5x_i^4$ (Lean names \lean{dF5_dx0}, \dots, \lean{dF5_dx5}); it is the formal partial derivative $\partial F_5/\partial x_i$, but differentiation itself is not formalized.

\section{Results}
\begin{theorem}[Euler identity; \lean{Hodge.VF5.euler_F5}]\label{thm:euler}
For all $x\in\mathbb{Z}^6$, $5F_5(x)=\sum_{i=0}^5 x_iD_i(x)$.
\end{theorem}

\begin{theorem}[Euler defect of $G$; \lean{Hodge.VF5.G_euler_defect}]\label{thm:defect}
For all $x\in\mathbb{Z}^6$,
\[
\sum_{i=0}^5 x_i\Bigl(5x_i^4-5\prod_{j\neq i}x_j\Bigr)-5G(x)=-5\prod_{j=0}^5x_j .
\]
\end{theorem}
The bracketed factors are the formal partials of $G$. Since $G$ is not homogeneous (its monomial $\prod x_j$ has degree~6), Euler's identity fails for it, and Theorem~\ref{thm:defect} computes the defect exactly.

\begin{lemma}\label{lem:prod}
\begin{enumerate}
\item (\lean{Hodge.VF5.partial_as_prod}) $5x^4=5\cdot x\cdot x\cdot x\cdot x$ for all $x\in\mathbb{Z}$.
\item (\lean{Hodge.VF5.int_five_mul_prod4_eq_zero}) If $x\in\mathbb{Z}$ and $5\cdot x\cdot x\cdot x\cdot x=0$, then $x=0$.
\end{enumerate}
\end{lemma}

\begin{theorem}[Isolated critical point of the affine cone; \lean{Hodge.VF5.affine_cone_isolated_at_origin}]\label{thm:iso}
Let $x\in\mathbb{Z}^6$. If $D_i(x)=0$ for all $i=0,\dots,5$, then $x_0=x_1=\dots=x_5=0$.
\end{theorem}
\begin{proof}
For each $i$, rewrite $D_i(x)=0$ as $5x_ix_ix_ix_i=0$ by Lemma~\ref{lem:prod}(1); since $\mathbb{Z}$ has no zero divisors, Lemma~\ref{lem:prod}(2) gives $x_i=0$.
\end{proof}

\begin{remark}
Theorem~\ref{thm:iso} is stated for integer points only. The analogous statement over $\mathbb{C}$ (which is what smoothness of $V(F_5)$ requires) follows by the same argument in any integral domain of characteristic $\neq5$, but that generalization is \emph{not} part of the formalization.
\end{remark}

\subsection*{Recorded constants and labels}
The following are checked by the kernel but are arithmetic or definitional facts, not geometric computations.
\begin{proposition}\label{prop:const}
\begin{enumerate}
\item (\lean{Hodge.VF5.h31_dim}) $126-6=120$ in $\mathbb{N}$.
\item (\lean{Hodge.VF5.sextic_h22_prim_container}) $1751=1751$ in $\mathbb{N}$.
\item (\lean{Hodge.VF5.sextic_Z_status_uncomputed}) The status string \lean{sextic_Z_status} equals \texttt{"uncomputed"}.
\item (\lean{Hodge.VF5.host_is_VF5}, \lean{Hodge.VF5.host_dim}) The documentation label has name \texttt{"V(F5)"} and dimension field $4$.
\end{enumerate}
\end{proposition}
Item~(1) is the subtraction one would use if $126$ and $6$ were the relevant dimensions in a Hodge-number calculation; item~(2) is a reflexivity statement recording a number. Neither the identification of these numbers with Hodge numbers, nor any cycle class, is proved here. Item~(3) records that a sextic quantity is explicitly marked as not computed.

\section{Scope}
Not formalized: schemes, projective space, smoothness as a geometric property, cohomology, Hodge numbers, primitive classes, algebraic cycles, and the Hodge conjecture. The paper claims exactly the integer identities and implications of Theorems~\ref{thm:euler}--\ref{thm:iso}, Lemma~\ref{lem:prod} and Proposition~\ref{prop:const}.

\section{Correspondence between claims and Lean declarations}
All declarations are in namespace \lean{Hodge.VF5}, file \lean{Hodge/VF5.lean}, Lean~4 core only.
\begin{center}
\begin{tabular}{@{}ll@{}}
\toprule
Claim & Lean declaration\\
\midrule
Theorem~\ref{thm:euler} & \lean{euler_F5}\\
Theorem~\ref{thm:defect} & \lean{G_euler_defect}\\
Lemma~\ref{lem:prod}(1) & \lean{partial_as_prod}\\
Lemma~\ref{lem:prod}(2) & \lean{int_five_mul_prod4_eq_zero}\\
Theorem~\ref{thm:iso} & \lean{affine_cone_isolated_at_origin}\\
Proposition~\ref{prop:const}(1) & \lean{h31_dim}\\
Proposition~\ref{prop:const}(2) & \lean{sextic_h22_prim_container}\\
Proposition~\ref{prop:const}(3) & \lean{sextic_Z_status_uncomputed}\\
Proposition~\ref{prop:const}(4) & \lean{host_is_VF5}, \lean{host_dim}\\
\bottomrule
\end{tabular}
\end{center}

\appendix
\section{Verification certificate}\label{app:verification-certificate}

The formal development was machine-checked as follows.

{\raggedright
\begin{itemize}
\item Repository:\\ {\footnotesize\url{https://github.com/BenFrohman/FermatQuinticHost}}
\item Verified commit:\\ {\footnotesize\texttt{0b56d4b210d9418931af559cc06fdda2ee5b9bd1}}
\item Lean toolchain: \texttt{leanprover/lean4:v4.22.0}
\item Mathlib: none (the development depends only on Lean core)
\item Continuous integration run (build, axiom audit, \texttt{sorry}/\texttt{axiom} scan; passed):\\ {\footnotesize\url{https://github.com/BenFrohman/FermatQuinticHost/actions/runs/41794839040}}
\item SHA-256 of the GitHub source tarball of the verified commit, downloaded twice with identical result:\\ {\footnotesize\url{https://github.com/BenFrohman/FermatQuinticHost/archive/0b56d4b210d9418931af559cc06fdda2ee5b9bd1.tar.gz}}\\
{\footnotesize\texttt{726065496f5fb2841642e615f7e7c0eb0fd7b3e6927d4bba7dfe4e2dd7d3568a}}
\end{itemize}\par}

The tarball hash refers to the verified commit rather than to a release tag, since a document cannot contain the hash of an archive that includes itself. The Lean sources are unchanged between that commit and the release tag.

\textbf{The development contains no \texttt{sorry} and no user-declared axioms.} Every one of the 10 theorems and lemmas listed below depends only on Lean's standard axioms \texttt{propext}, \texttt{Classical.choice} and \texttt{Quot.sound} (or on none). The table is the verbatim output of \texttt{lake env lean AxiomAudit.lean} (\texttt{\#print axioms}) on the verified commit.

{\footnotesize
\begin{longtable}{p{0.55\linewidth}p{0.38\linewidth}}
\hline
Lean declaration & \texttt{\#print axioms} \\
\hline\endhead
\path{Hodge.VF5.G_euler_defect} & \texttt{[propext, Classical.choice, Quot.sound]} \\
\path{Hodge.VF5.affine_cone_isolated_at_origin} & \texttt{[propext, Classical.choice, Quot.sound]} \\
\path{Hodge.VF5.euler_F5} & \texttt{[propext, Classical.choice, Quot.sound]} \\
\path{Hodge.VF5.h31_dim} & \texttt{(no axioms)} \\
\path{Hodge.VF5.host_dim} & \texttt{(no axioms)} \\
\path{Hodge.VF5.host_is_VF5} & \texttt{(no axioms)} \\
\path{Hodge.VF5.int_five_mul_prod4_eq_zero} & \texttt{[propext]} \\
\path{Hodge.VF5.partial_as_prod} & \texttt{[propext, Classical.choice, Quot.sound]} \\
\path{Hodge.VF5.sextic_Z_status_uncomputed} & \texttt{(no axioms)} \\
\path{Hodge.VF5.sextic_h22_prim_container} & \texttt{(no axioms)} \\
\hline
\end{longtable}}

\begin{thebibliography}{9}
\bibitem{shioda} T.~Shioda, \emph{The Hodge conjecture for Fermat varieties}, Math.\ Ann.\ \textbf{245} (1979), 175--184.
\bibitem{lean} L.~de~Moura and S.~Ullrich, \emph{The Lean~4 theorem prover and programming language}, CADE-28, LNCS 12699 (2021), 625--635.
\end{thebibliography}
\end{document}
